\documentclass{beamer}

\mode<presentation>
\title{SoC1\&2: Physical Verification: DRC, LVS, ERC}
\author[AS]{Andreas Siggelkow \newline \texttt{andreas.siggelkow@rwu.de}}
\date[2026]{\today}
\subtitle{RWU}
\institute{Elektrotechnik und Informatik}

%%%%%%%%%%%%%%%%%%%%%%%%%%%%%%%%
% Logo
%%%%%%%%%%%%%%%%%%%%%%%%%%%%%%%%
\pgfdeclareimage[height=1cm]{hslogo}{./hs}
\logo{\pgfuseimage{hslogo}}
%%%%%%%%%%%%%%%%%%%%%%%%%%%%%%%%
% Macros
%%%%%%%%%%%%%%%%%%%%%%%%%%%%%%%%
\usepackage[T1]{fontenc}
\usepackage{lmodern}
\usepackage{textcomp}
\definecolor{links}{HTML}{2A1B81}
\hypersetup{colorlinks,linkcolor=,urlcolor=links}
\usepackage{graphicx}
\usepackage{ragged2e, tabularx}
\newcolumntype{Y}{>{\RaggedRight\arraybackslash}X}
\usetheme{Singapore}
\usecolortheme{seahorse}
\usefonttheme[onlylarge]{structuresmallcapsserif}
\usefonttheme[onlysmall]{structurebold}
\setbeamertemplate{footline}[frame number]
\usepackage{listings}
\newcommand{\astitle}{nix}
\newcommand{\assubtitle}{nix}
\newcommand{\asfolie}[3]{\begin{frame}\frametitle{#1}\begin{block}{#2}{#3}\end{block}\end{frame}}

% Commands and tool output are set with listings: dashes, $, _, \ and quotes
% appear exactly as they have to be typed.
\definecolor{lstbg}{HTML}{F2F2F7}
\definecolor{lstcomment}{HTML}{5A5A8A}
\lstset{
  basicstyle=\ttfamily\scriptsize,
  columns=fullflexible,
  keepspaces=true,
  upquote=true,
  breaklines=true,
  breakatwhitespace=false,
  backgroundcolor=\color{lstbg},
  frame=none,
  xleftmargin=2pt,
  xrightmargin=2pt,
  aboveskip=4pt,
  belowskip=4pt,
  showstringspaces=false
}
\lstdefinestyle{sh}{
  morecomment=[l]{\#},
  commentstyle=\color{lstcomment}\itshape
}
\lstdefinestyle{out}{basicstyle=\ttfamily\tiny}
\newcommand{\cmd}[1]{\texttt{#1}}





\begin{document}
\setcounter{tocdepth}{1}

\begin{frame}
  \titlepage
\end{frame}

\begin{frame}
  \frametitle{Outline}
  \tableofcontents
\end{frame}

%%%%%%%%%%%%%%%%%%%%%%%%%%%%%%%%
% Overview
%%%%%%%%%%%%%%%%%%%%%%%%%%%%%%%%
\section{Overview}
\renewcommand{\astitle}{Overview}

\renewcommand{\assubtitle}{The Last Checks Before Tape-out}
\begin{frame}
\frametitle{\astitle}
\begin{block}{\assubtitle}
  \footnotesize
  \begin{tabularx}{\linewidth}{@{}lY@{}}
    \textbf{check} & \textbf{question}\\\hline
    DRC & does the layout obey the manufacturing rules of XH018: widths,
          spacings, enclosures, densities, antenna ratios?\\
    LVS & is the layout the circuit of the netlist: same devices, same
          connections, same ports?\\
    ERC & electrical rules: floating gates, wells not tied to a supply,
          shorts between supplies\\
  \end{tabularx}
  \vspace{4pt}
  \begin{itemize}
    \item Input: the \emph{merged GDS}, the full layout incl.\ all cells,
      not the LEF abstracts Innovus works with
    \item Rules: the X-FAB rule decks (runsets) -- the fab only accepts a
      GDS that passes \emph{its} decks
    \item Innovus \cmd{verify\_drc} only knows the LEF rules of the routing
      layers: a hint, not a sign-off
  \end{itemize}
\end{block}
\end{frame}

\renewcommand{\assubtitle}{Which Tool?}
\begin{frame}
\frametitle{\astitle}
\begin{block}{\assubtitle}
  \begin{itemize}
    \item The PDK decides: X-FAB ships rule decks for certain tools
    \begin{itemize}
      \item \textbf{Pegasus} (Cadence, PVL decks; also reads PVS decks)
      \item \textbf{Calibre} (Siemens, SVRF decks)
      \item Assura (Cadence, older, inside Virtuoso)
    \end{itemize}
    \item The lab runs Cadence tools: Pegasus first. Calibre needs its
      own licence
    \item The scripts support both: \cmd{PV\_TOOL} in \cmd{\$XFAB\_PATHS}
      or \cmd{pv.sh -t calibre}
    \item Find out in the lab which decks the PDK has:
      \cmd{ls} in the PDK tree, directories named after the tools
  \end{itemize}
\end{block}
\end{frame}

\renewcommand{\assubtitle}{Where in the Flow?}
\begin{frame}
\frametitle{\astitle}
\begin{block}{\assubtitle}
  \footnotesize
  \begin{tabularx}{\linewidth}{@{}lY@{}}
    \textbf{step} & \textbf{result}\\\hline
    \cmd{innovus\_post.sh} (05) & \cmd{cc\_chip\_top.gds} (merged, with pad
                                labels), \cmd{cc\_chip\_top\_lvs.v}\\
    \textbf{DRC}              & result count per deck: must be 0\\
    \textbf{LVS + ERC}        & CORRECT, 0 ERC results\\
    fix                       & in Innovus (or the RTL), new GDS, again\\
  \end{tabularx}
  \vspace{4pt}
  \begin{itemize}
    \item STA, LEC and GLS check the logic and the timing; DRC/LVS check
      the geometry. All of them must pass on the \emph{same} layout
  \end{itemize}
\end{block}
\end{frame}

%%%%%%%%%%%%%%%%%%%%%%%%%%%%%%%%
% Inputs
%%%%%%%%%%%%%%%%%%%%%%%%%%%%%%%%
\section{Inputs}
\renewcommand{\astitle}{Inputs}

\renewcommand{\assubtitle}{From Innovus}
\begin{frame}
\frametitle{\astitle}
\begin{block}{\assubtitle}
  \begin{itemize}
    \item \cmd{cc\_chip\_top.gds} (\cmd{02c\_rv64\_gds.tcl}):
      \cmd{write\_stream -merge} with the GDS of standard cells, IO cells
      and SPRAM macros; \cmd{cc\_chip\_top\_gds.rpt} must show
      \cmd{missing: 0}
    \item Pad labels: one text per pad at its centre on
      \cmd{LABEL\_LAYER}: the port name, or \cmd{VDD}/\cmd{GND}/\cmd{VDDO}/
      \cmd{GNDO} on supply pads (list: \cmd{cc\_chip\_top\_labels.txt}).
      LVS uses them to match ports and supply nets
    \item \cmd{cc\_chip\_top\_lvs.v} (\cmd{02b}): netlist \emph{with} supply
      ports and physical-only cells (fillers, corners) -- LVS compares
      device by device
  \end{itemize}
\end{block}
\end{frame}

\renewcommand{\assubtitle}{From the PDK: \texttt{\$XFAB\_PATHS}}
\begin{frame}[fragile]
\frametitle{\astitle}
\begin{block}{\assubtitle}
\begin{lstlisting}[style=sh]
set PV_TOOL     pegasus                  ;# pegasus | calibre
set DRC_RULES   [list <drc_deck> <antenna_deck>]   ;# one run each
set LVS_RULES   <lvs_deck>               ;# incl. ERC checks
set PV_SWITCHES [list]                   ;# #DEFINE switches of the decks
set CDL_FILES   [list <D_CELLS_HD.cdl> <IO_CELLS_3V.cdl> \
                      <imem.cdl> <dmem.cdl>]
set LABEL_LAYER <top_metal>              ;# pad labels (02c)
\end{lstlisting}
  \begin{itemize}
    \item \cmd{PV\_SWITCHES}: the decks cover several process options; the
      switches select ours (4 metals, LPMOS, top-metal thickness). The
      names are in the header of each deck
    \item \cmd{CDL\_FILES}: transistor netlists of the cells; the SPRAM
      CDL comes with the memory compiler
  \end{itemize}
\end{block}
\end{frame}

%%%%%%%%%%%%%%%%%%%%%%%%%%%%%%%%
% Scripts
%%%%%%%%%%%%%%%%%%%%%%%%%%%%%%%%
\section{Scripts}
\renewcommand{\astitle}{Scripts}

\renewcommand{\assubtitle}{Control Files Instead of Edited Decks}
\begin{frame}
\frametitle{\astitle}
\begin{block}{\assubtitle}
  \begin{itemize}
    \item The X-FAB decks are never changed (NDA, and a changed deck is no
      longer the fab's sign-off)
    \item \cmd{05\_rv64\_pv.tcl} (tclsh) writes small control files into
      \cmd{build\_pv/}: inputs, report names, switches, then
      \cmd{include} of the deck
    \begin{itemize}
      \item \cmd{drc\_1.ctl}, \cmd{drc\_2.ctl}, \dots: one per deck in
        \cmd{DRC\_RULES}
      \item \cmd{lvs.ctl}: LVS with ERC
    \end{itemize}
    \item Same idea as \cmd{quantus.ccl}: the build directory holds the
      X-FAB paths, the repository only the generator
  \end{itemize}
\end{block}
\end{frame}

\renewcommand{\assubtitle}{\texttt{lvs.ctl} for Pegasus}
\begin{frame}[fragile]
\frametitle{\astitle}
\begin{block}{\assubtitle}
\begin{lstlisting}[style=sh]
layout_path "<...>/cc_chip_top.gds" gdsii;
layout_primary "cc_chip_top";
schematic_path "<D_CELLS_HD.cdl>" cdl;     // one line per CDL file
schematic_path "<...>/cc_chip_top_lvs.v" verilog;
schematic_primary "cc_chip_top";
report_lvs "lvs.report" -replace;
results_db -erc "erc.results" -ascii;
report_summary -erc "erc.summary" -replace;
lvs_power_name "VDD" "VDDO";
lvs_ground_name "GND" "GNDO";
#DEFINE <switch>
include "<lvs_deck>";
\end{lstlisting}
  \begin{itemize}
    \item Pegasus reads the Verilog netlist directly, the cells as CDL
  \end{itemize}
\end{block}
\end{frame}

\renewcommand{\assubtitle}{\texttt{lvs.ctl} for Calibre}
\begin{frame}[fragile]
\frametitle{\astitle}
\begin{block}{\assubtitle}
\begin{lstlisting}[style=sh]
LAYOUT PATH "<...>/cc_chip_top.gds"
LAYOUT PRIMARY "cc_chip_top"
LAYOUT SYSTEM GDSII
SOURCE PATH "cc_chip_top_lvs.sp"        // from v2lvs
SOURCE PRIMARY "cc_chip_top"
SOURCE SYSTEM SPICE
LVS REPORT "lvs.report"
ERC RESULTS DATABASE "erc.results" ASCII
LVS POWER NAME "VDD" "VDDO"
LVS GROUND NAME "GND" "GNDO"
#DEFINE <switch>
INCLUDE "<lvs_deck>"
\end{lstlisting}
  \begin{itemize}
    \item Calibre wants SPICE: \cmd{pv.sh} first runs
      \cmd{v2lvs -v cc\_chip\_top\_lvs.v -s <cdl> \dots\ -s0 GND -s1 VDD}
  \end{itemize}
\end{block}
\end{frame}

\renewcommand{\assubtitle}{Running}
\begin{frame}[fragile]
\frametitle{\astitle}
\begin{block}{\assubtitle}
\begin{lstlisting}[style=sh]
cd RWU_RV/RV_NoPipeline
./scripts/innovus/innovus_post.sh -b    # 05, up to stage gds
./scripts/pv/pv.sh                      # DRC + LVS/ERC, PV_TOOL
./scripts/pv/pv.sh -m drc -j 8          # DRC only, 8 CPUs
./scripts/pv/pv.sh -m lvs -t calibre    # LVS with Calibre
\end{lstlisting}
\begin{lstlisting}[style=out]
=== [pv] DRC drc_1 (pegasus)
    <drc_deck>: N DRC result(s)
=== [pv] DRC drc_2 (pegasus)
    <antenna_deck>: N DRC result(s)
=== [pv] LVS + ERC (pegasus)
    LVS: CORRECT
    ERC: 0 result(s) (erc.summary)
PASS  physical verification of cc_chip_top (all, pegasus)
\end{lstlisting}
\end{block}
\end{frame}

%%%%%%%%%%%%%%%%%%%%%%%%%%%%%%%%
% DRC
%%%%%%%%%%%%%%%%%%%%%%%%%%%%%%%%
\section{DRC}
\renewcommand{\astitle}{DRC}

\renewcommand{\assubtitle}{Reading the Results}
\begin{frame}
\frametitle{\astitle}
\begin{block}{\assubtitle}
  \begin{itemize}
    \item \cmd{drc\_<n>.summary}: count per rule. Sort by count: a few rules
      usually cause most results
    \item \cmd{drc\_<n>.results}: every result with its coordinates
    \item Look at them on the layout:
    \begin{itemize}
      \item Calibre: \cmd{calibre -rve drc\_1.results} (with Virtuoso or
        DESIGNrev)
      \item Pegasus: Pegasus DesignReview, or load the results into
        Innovus as markers (violation browser) -- VERIFY-LAB
    \end{itemize}
    \item Many results in one rule, all at the macros or the pads? Then
      it is a cell or a merge problem, not the routing
  \end{itemize}
\end{block}
\end{frame}

\renewcommand{\assubtitle}{Typical First-Run Results}
\begin{frame}
\frametitle{\astitle}
\begin{block}{\assubtitle}
  \footnotesize
  \begin{tabularx}{\linewidth}{@{}lY@{}}
    \textbf{rule type} & \textbf{fix}\\\hline
    metal density    & fill in Innovus (\cmd{FILL\_LAYERS}) too sparse,
                       or X-FAB wants its own fill deck\\
    antenna          & left over after the router's fixing: more diodes,
                       check \cmd{ANTENNA\_CELL}, LEF antenna data\\
    well/implant gaps & missing core fillers or well taps
                       (\cmd{CORE\_FILLERS}, \cmd{WELLTAP\_CELL})\\
    at the macros    & stripes over the SPRAM: blocked layers, or halo
                       too small\\
    at the pads      & IO fillers/corners, pad spacing (TECH-06)\\
    empty cells      & \cmd{cc\_chip\_top\_gds.rpt}: a GDS file missing in
                       the merge\\
  \end{tabularx}
\end{block}
\end{frame}

%%%%%%%%%%%%%%%%%%%%%%%%%%%%%%%%
% LVS
%%%%%%%%%%%%%%%%%%%%%%%%%%%%%%%%
\section{LVS and ERC}
\renewcommand{\astitle}{LVS and ERC}

\renewcommand{\assubtitle}{How LVS Works}
\begin{frame}
\frametitle{\astitle}
\begin{block}{\assubtitle}
  \begin{itemize}
    \item \textbf{Extraction}: from the layout geometry LVS recognizes
      the devices (transistors, diodes, resistors) and the nets connecting
      them
    \item \textbf{Source}: the netlist is expanded down to the same
      devices via the cell CDL
    \item \textbf{Compare}: both graphs are matched, starting from the
      labelled nets (ports, supplies)
    \item Result \cmd{CORRECT}: same devices, same connections, same
      ports. Otherwise a list of mismatches in \cmd{lvs.report}
  \end{itemize}
\end{block}
\end{frame}

\renewcommand{\assubtitle}{Typical Mismatches}
\begin{frame}
\frametitle{\astitle}
\begin{block}{\assubtitle}
  \footnotesize
  \begin{tabularx}{\linewidth}{@{}lY@{}}
    \textbf{symptom} & \textbf{cause}\\\hline
    ports not found      & pad labels missing or on the wrong layer
                           (\cmd{LABEL\_LAYER}, deck's text layer)\\
    VDD and GND merged   & a short in the power grid -- first check
                           \cmd{cc\_chip\_top\_pg\_connectivity.rpt}\\
    open supply net      & a row, macro or pad not connected to the ring\\
    extra devices        & physical-only cells missing in the LVS netlist
                           (\cmd{-include\_phys\_insts})\\
    unknown cell         & CDL of a cell missing in \cmd{CDL\_FILES}\\
    VDDO vs.\ VDD        & IO cells: core and IO supply pins in the wrong
                           \cmd{PG\_PINS\_*} list\\
  \end{tabularx}
  \vspace{4pt}
  \begin{itemize}
    \item ERC (\cmd{erc.summary}) often points to the same root cause as
      an LVS mismatch: fix LVS first
  \end{itemize}
\end{block}
\end{frame}

%%%%%%%%%%%%%%%%%%%%%%%%%%%%%%%%
% Open points
%%%%%%%%%%%%%%%%%%%%%%%%%%%%%%%%
\section{Open Points}
\renewcommand{\astitle}{Open Points}

\renewcommand{\assubtitle}{VERIFY-LAB: Not Yet Run}
\begin{frame}
\frametitle{\astitle}
\begin{block}{\assubtitle}
  \begin{itemize}
    \item Which decks the PDK has, and their switches
    \item Whether the X-FAB decks already set layout/source paths
      themselves (then use the deck's variables instead of the control
      file statements)
    \item PVL statement names for Pegasus, the \cmd{pegasus}/\cmd{calibre}
      command lines
    \item Summary and report wording that \cmd{pv.sh} looks for
    \item Pad labels: \cmd{add\_text} in Innovus, the text layer in the
      stream map, the text layer the LVS deck expects
  \end{itemize}
\end{block}
\end{frame}

\renewcommand{\assubtitle}{X-FAB NDA}
\begin{frame}
\frametitle{\astitle}
\begin{block}{\assubtitle}
  \begin{itemize}
    \item GDS, CDL, rule decks, control files, reports: all contain X-FAB
      data or names
    \item \cmd{build\_innovus/} and \cmd{build\_pv/} are in
      \cmd{.gitignore}: never commit them, never copy them out of the lab
    \item The GDS goes to X-FAB only through the official tape-out
      channel
  \end{itemize}
\end{block}
\end{frame}

\renewcommand{\assubtitle}{Before Tape-out}
\begin{frame}
\frametitle{\astitle}
\begin{block}{\assubtitle}
  \begin{itemize}
    \item DRC 0, LVS CORRECT, ERC 0 -- on the final GDS
    \item STA, LEC and the post-route GLS on the same layout
    \item Pad order and supply pads checked (TECH-06), seal ring and
      logo as X-FAB requires
    \item X-FAB's own tape-out checks (their DRC run on the submitted GDS)
  \end{itemize}
\end{block}
\end{frame}

\section{Q\&A}
\renewcommand{\astitle}{Q\&A}
\begin{frame}
\frametitle{\astitle}
\vfill
\centering
\begin{Huge}\textcolor{red}{Questions?}\end{Huge}
\vfill
\end{frame}

\end{document}
